\documentclass[12pt]{article}

% === Packages ===
\usepackage[margin=1in]{geometry}
\usepackage{amsmath, amssymb}
\usepackage{setspace}
\usepackage{titlesec}
\usepackage{enumitem}
\usepackage{helvet}
\renewcommand{\familydefault}{\sfdefault} % modern sans-serif look

% === Formatting tweaks ===
\setstretch{1.15}
\setlength{\emergencystretch}{3em}
\setlist[enumerate]{itemsep=0.4em, topsep=0.2em}
\titleformat{\section}{\large\bfseries\sffamily}{\thesection.}{1em}{}
\titleformat{\subsection}{\large\bfseries\sffamily}{\thesubsection}{1em}{}

% === Title ===
\title{\vspace{-2em}\textbf{Applied Mathematics Oral Exam Questions and Model Answers}\\
\large Immersed Boundary Method}
\author{Prepared for Graduate Oral Examination}
\date{}

\begin{document}
\maketitle

Unless stated otherwise, the answers describe the classical massless immersed
boundary formulation: the structure is advected by the interpolated fluid
velocity, and its force is measured with respect to a stated Lagrangian
reference measure.

\subsection*{Conceptual Foundations}

\begin{enumerate}
\item \textbf{What is the immersed boundary (IB) method, and what class of problems was it originally designed to study?}

\textit{Sample answer:}
The immersed boundary method is a computational framework for fluid--structure
interaction (FSI). Charles Peskin developed it to model cardiac fluid dynamics,
including the interaction of flexible heart structures with viscous blood. Its
defining feature is an Eulerian description of the fluid coupled to a
Lagrangian description of the immersed structure without requiring the fluid
grid to conform to the moving interface.

\item \textbf{Describe the Eulerian--Lagrangian coupling used in the immersed boundary method.}

\textit{Sample answer:}
The fluid velocity and pressure live on a fixed Eulerian grid, while material
points or finite elements describe the structure. A discrete delta kernel
spreads Lagrangian forces to the fluid grid, and the adjoint interpolation
operator evaluates fluid velocity at the structure. In the classical massless
formulation, that interpolated velocity advances the structure. The integration
weights and the units of the Lagrangian force must match the chosen reference
measure.

\item \textbf{What does it mean that boundary conditions are enforced “weakly” in the immersed boundary method?}

\textit{Sample answer:}
Here ``weakly'' means that the kinematic constraint is represented through
regularized interpolation and force coupling rather than imposed as a sharp
pointwise condition on a body-fitted fluid grid. The interface has a numerical
thickness of a few grid cells, so pointwise traction jumps are smoothed at finite
resolution. In the singular continuum formulation, a no-slip interface normally
has continuous velocity, while pressure and/or viscous traction can jump when
the interface carries the corresponding singular force; such jumps are not
automatic for every immersed structure.
\end{enumerate}

\subsection*{Mathematical and Numerical Structure}

\begin{enumerate}[resume]
\item \textbf{Write down the generic form of the force coupling in the immersed boundary method.}

\textit{Sample answer:}
Let $\mathbf{q}$ label the reference structure, let
$\mathbf{X}(\mathbf{q},t)$ be its current position, and let
$\mathbf{F}(\mathbf{q},t)$ be force density with respect to the chosen
reference measure $d\mathbf{q}$. Then, with the exact Dirac delta,
\[
\mathbf{f}(\mathbf{x},t) = \int_U \mathbf{F}(\mathbf{q},t)\,
\delta\!\left(\mathbf{x}-\mathbf{X}(\mathbf{q},t)\right)\,d\mathbf{q},
\]
and, in the classical massless formulation,
\[
\frac{\partial\mathbf{X}}{\partial t}(\mathbf{q},t)
= \int_\Omega \mathbf{u}(\mathbf{x},t)\,
\delta\!\left(\mathbf{x}-\mathbf{X}(\mathbf{q},t)\right)\,d\mathbf{x}.
\]
At the discrete level, quadrature factors are essential and spreading and
interpolation should be adjoint under the chosen inner products.

\item \textbf{What factors control the accuracy of an immersed boundary simulation?}

\textit{Sample answer:}
Accuracy depends on the Eulerian discretization, Lagrangian quadrature and mesh
quality, relative mesh spacings, kernel support and moment properties, time
integration, and solver tolerances. A Lagrangian mesh that is too coarse can
leave gaps; extreme oversampling can worsen conditioning without adding
independent information. Refinement can converge to the intended continuum
problem only when the complete coupled scheme is consistent and stable; a
regularized kernel alone does not guarantee a rate or even convergence.

\item \textbf{Why can stiff elastic forces lead to severe time-step restrictions in explicit IB schemes?}

\textit{Sample answer:}
Large elastic eigenvalues coupled through the fluid create fast modes. An
explicit update must resolve those modes, and its stable time step can shrink
rapidly with mesh refinement and stiffness. Semi-implicit, fully implicit, or
split preconditioned methods can relax this restriction, although they introduce
linear or nonlinear solves and do not remove accuracy requirements.
\end{enumerate}

\subsection*{Advantages, Limitations, and Failure Modes}

\begin{enumerate}[resume]
\item \textbf{List two major advantages and two major disadvantages of the immersed boundary method.}

\textit{Sample answer:}
Advantages include handling complex moving geometry and large deformation on a
fixed fluid grid, and a common coupling framework for fibers, surfaces, and
volumetric solids. Disadvantages include regularization error across a
grid-scale interface, limited sharp-interface accuracy for pressure or traction,
possible volume leakage if the coupled meshes and kernel are poorly resolved,
and severe explicit time-step restrictions for stiff structures.

\item \textbf{Give an example of a common failure mode in IB simulations and explain its cause.}

\textit{Sample answer:}
A common failure mode is leakage through an interface whose Lagrangian spacing
is too coarse relative to the Eulerian grid and kernel support, leaving the
regularized force layers unable to form an effective barrier. Conversely,
extreme Lagrangian oversampling can create redundant constraints or poor
conditioning. Severe time-step restrictions primarily arise from the physical
or penalty stiffness together with explicit coupling; fine spacing can amplify
the largest discrete elastic eigenvalues but is not itself a material stiffness.
\end{enumerate}

\subsection*{Comparison with Other FSI Methods}

\begin{enumerate}[resume]
\item \textbf{Compare the immersed boundary method with grid-fitted (ALE or body-fitted) FSI methods.}

\textit{Sample answer:}
Body-fitted/ALE methods align the fluid mesh with the interface, enabling sharp
boundary conditions and accurate tractions, but mesh motion can become difficult
under large deformation or topological change and may require remeshing. IB
methods retain a fixed, nonconforming fluid grid and simplify large motion at the
cost of regularization error and more delicate interface-resolution and
preconditioning issues. Neither family is uniformly more accurate or robust;
the tradeoff depends on geometry, discretization, and quantity of interest.

\item \textbf{How does the immersed interface method (IIM) differ from the immersed boundary method?}

\textit{Sample answer:}
Immersed interface methods incorporate analytically derived jump conditions into
stencils that cross an interface, giving a sharp representation and potentially
higher local accuracy for sufficiently smooth, well-tracked interfaces. The
classical IB method represents singular coupling with a regularized delta kernel,
giving a diffuse numerical interface. The comparison is not simply ``exact''
versus ``inexact'': both have discretization error, and their applicability
depends on whether the jump data and interface geometry can be supplied and
updated reliably.
\end{enumerate}

\newpage
\subsection*{Applications and Practical Considerations}

\begin{enumerate}[resume]
\item \textbf{Give two biological and two engineering applications of the immersed boundary method.}

\textit{Sample answer:}
Biological applications include cardiac flow and locomotion of organisms such
as fish or jellyfish. Engineering applications include flow-induced vibration
of flexible beams or flags and fluid interaction with soft robotic components.
For each example one must decide whether a fiber, shell, membrane, or volumetric
solid model supplies the appropriate mechanics.

\item \textbf{Why is the availability of open-source IB software an important practical advantage?}

\textit{Sample answer:}
Open-source frameworks such as IBAMR and IB2d expose algorithms and examples,
lower implementation barriers, and support reproducibility. They do not replace
verification, convergence studies, or validation, and results may depend on
version, solver tolerances, kernel, and discretization choices. Recording those
choices is especially important for large three-dimensional calculations.
\end{enumerate}

\subsection*{References}

The model answers are adapted from the Module~1 decks (1-1 and 1-2), which
follow these sources.
\begin{itemize}
\item C.~S. Peskin, ``Numerical analysis of blood flow in the heart,''
  \emph{J. Comput. Phys.} 25 (1977), 220--252,
  doi:10.1016/0021-9991(77)90100-0. Origin of the method in cardiac flow
  (Question~1).
\item C.~S. Peskin, ``The immersed boundary method,'' \emph{Acta Numerica} 11
  (2002), 479--517, doi:10.1017/S0962492902000077. Coupling equations,
  adjointness, kernels, and stiffness.
\item B.~E. Griffith and N.~A. Patankar, ``Immersed methods for
  fluid--structure interaction,'' \emph{Annu. Rev. Fluid Mech.} 52 (2020),
  421--448, doi:10.1146/annurev-fluid-010719-060228. Comparison with
  body-fitted, ALE, and sharp-interface methods.
\item Z. Li and M.-C. Lai, ``The immersed interface method for the
  Navier--Stokes equations with singular forces,'' \emph{J. Comput. Phys.} 171
  (2001), 822--842. Immersed interface method (Question~10).
\item IBAMR, \texttt{https://ibamr.github.io}; N.~A. Battista, W.~C. Strickland,
  and L.~A. Miller, ``IB2d: a Python and MATLAB implementation of the immersed
  boundary method,'' \emph{Bioinspir. Biomim.} 12 (2017), 036003
  (Question~12).
\end{itemize}

\end{document}
